\section{Sixteen-qubit conditional switching construction}
\label{app:switching16}

\subsection{Cycle basis and reference state}

We specialize the binary-linear encoder construction to the cycle space of the graph in Section~\ref{sec:conditional-switching16}. We use column vectors throughout, with all matrix operations over $\mathbb F_2$.

We order the sixteen edges lexicographically as
\begin{equation}
 (01,02,03,04,05,06,12,13,14,15,16,23,24,25,26,34).
\end{equation}
The six root edges $01,\ldots,06$ form a spanning tree. The remaining edges form the chord set $\mathcal C=\{12,13,14,15,16,23,24,25,26,34\}$. Each chord and its two root edges define a fundamental triangle.

We represent these ten triangles by a matrix $B\in\mathbb F_2^{10\times16}$. Each row contains ones on its triangle's three edges and zeros elsewhere. The chord columns form the identity matrix, so $B$ has rank ten. The connected graph has cycle rank $16-7+1=10$, and these rows therefore form a basis of its cycle space. The encoding matrix in Section~\ref{sec:input-symmetry} is $C=B^{\mathsf T}$.

The six independent vertex switches generate the group $G$ of $64$ switching masks. Each switch selects either zero or two edges of every triangle, so $Bg=0$ for every $g\in G$. Since $\ker B$ has dimension $16-10=6$, it also contains $64$ masks. Thus, $G=\ker B$.

For a vector of chord bits $r\in\mathbb F_2^{10}$, the edge configuration is $B^{\mathsf T}r$. We define the reference state as
\begin{align}
 \ket{\psi_\beta}
 &=\sum_{r\in\mathbb F_2^{10}}a_\beta(r)\ket{B^{\mathsf T}r},\\
 a_\beta(r)&=c_\beta^{10-|r|}(-i s_\beta)^{|r|},
\end{align}
where $|r|$ is the Hamming weight of $r$,
\begin{equation}
 c_\beta=\cos(t_\beta/2),\qquad
 s_\beta=\sin(t_\beta/2),\qquad
 t_\beta=\frac{\pi}{2}\tanh(1.5\beta).
\end{equation}
The normalized rotation factors and injectivity of $B^{\mathsf T}$ ensure normalization, as established in Appendix~\ref{app:binary-linear-encoders}. Applying the Hadamard layer gives the encoded input $\ket{\phi_\beta}=H^{\otimes16}\ket{\psi_\beta}$.

\subsection{Classical sampling and amplitude ratios}

We sample the reference distribution by drawing the ten chord bits independently, with probability $s_\beta^2$ of obtaining one, and returning $B^{\mathsf T}r$. The chord entries of this configuration equal $r$ because the chord columns of $B$ form the identity matrix. They therefore also allow us to recover the latent bits from any reference sample.

For a bit-flip mask $m$, we collect its ten chord entries into $u$ and test whether $B^{\mathsf T}u=m$. If the test fails, $m$ lies outside the cycle space, and flipping a supported reference configuration by $m$ gives zero amplitude. If the test succeeds, the flip maps $B^{\mathsf T}r$ to $B^{\mathsf T}(r\oplus u)$.

For $\beta\neq0$, both $c_\beta$ and $s_\beta$ are nonzero, so every cycle-space configuration has nonzero reference amplitude. The ratio for a mask within the cycle space is
\begin{equation}
 \frac{a_\beta(r\oplus u)}{a_\beta(r)}
 =
 \left[-i\tan\!\left(\frac{t_\beta}{2}\right)\right]^{|r\oplus u|-|r|}.
\end{equation}
Masks outside the cycle space give zero ratio. At $\beta=0$, the reference state is $\ket{0}^{\otimes16}$, so the ratio equals one for $m=0$ and zero otherwise. We evaluate ratios only at sampled configurations with nonzero reference amplitude. This procedure provides classical sampling and amplitude-ratio access without enumerating the full sixteen-qubit distribution.

\subsection{Switching symmetry}

For every switching mask $g\in G$, the relation $Bg=0$ gives
\begin{equation}
 g^{\mathsf T}B^{\mathsf T}r=(Bg)^{\mathsf T}r=0.
\end{equation}
Every reference configuration therefore has even parity on the edges selected by $g$. Consequently,
\begin{equation}
 Z_g\ket{\psi_\beta}=\ket{\psi_\beta},
 \qquad
 X_g\ket{\phi_\beta}=\ket{\phi_\beta}.
\end{equation}
The IQP unitary commutes with $X_g$, so applying it preserves this eigenvalue. The output amplitudes at $x$ and $x\oplus g$ are therefore equal, giving
\begin{equation}
 q_{\boldsymbol\theta}(x\mid\beta)
 =
 q_{\boldsymbol\theta}(x\oplus g\mid\beta),
 \qquad g\in G.
\end{equation}
This equality holds for every condition and IQP parameter vector. The encoder thus enforces the six independent vertex-switching symmetries, independently of training accuracy.
